Анализ графиков показывает:
\begin{itemize}
\item При малой интенсивности ($\lambda \lesssim 0{,}05$ 1/с) производительности систем практически совпадают
и равны $\lambda$, так как потери малы. С ростом $\lambda$ производительность СИСТЕМЫ\_2 быстро
насыщается на уровне $1/b = 0{,}1$ 1/с, а СИСТЕМЫ\_1 продолжает расти к пределу $0{,}2$ 1/с.
Таким образом, при любом $\lambda$ выполняется $\Lambda'_1 \ge \Lambda'_2$.
\item Вероятность потери в СИСТЕМЕ\_2 меньше только при $\lambda < 0{,}054$ 1/с (там общий накопитель на 3 места
эффективнее раздельных по 2), при большей интенсивности потери в СИСТЕМЕ\_1 ниже.
\item Загрузка прибора СИСТЕМЫ\_2 всегда выше: один прибор принимает весь поток, а в СИСТЕМЕ\_1
поток делится на два прибора.
\item Время пребывания в СИСТЕМЕ\_1 меньше при всех $\lambda$, разница растёт с нагрузкой.
\item Суммарная длина очередей и число заявок в СИСТЕМЕ\_1 больше, чем в СИСТЕМЕ\_2, при
$\lambda > 0{,}47$ и $\lambda > 0{,}26$ 1/с соответственно, так как суммарная ёмкость СИСТЕМЫ\_1 больше
(2 прибора + 4 места против 1 прибора + 3 мест). При заданном $\lambda=0{,}5$ это проявляется
незначительно превышением $L$.
\end{itemize}
